\documentclass[../main]{subfiles}
\begin{document}
\chapter{Localización}
\img{images/12/mapa.png}{Mapa de Localización.}

\info{Estudio localización}{El estudio de localización tiene como propósito encontrar la ubicación más ventajosa para el proyecto, es decir, la opción que, cubriendo las exigencias o requerimientos del proyecto, contribuya a minimizar los costos de inversión y los costos y gastos durante el período productivo.}

\subsection{Estudio de localización}



Al considerar la localización del proyecto es posible concluir que hay más de una
solución factible adecuada por lo que hay que determinar las variables relevantes
más importantes en forma concluyente. De igual forma, una localización que se ha
determinado como óptima en las condiciones vigentes puede no serlo en el futuro.
Por lo tanto, la selección de la ubicación debe considerar su carácter definitivo o
transitorio.

\img{images/12/macromicro.png}{Pasos para el estudio de localización}

Sin embargo, el problema tampoco es puramente económico. Los factores
técnicos, legales, tributarios, sociales, etc. (figura) deben tomarse en
consideración, aunque en algunos casos, los efectos en el resultado del proyecto no pueden ser considerados en factores monetarios, por lo que siempre quedará una variable subjetiva que pudiera afectar la decisión, como por ejemplo, intereses
personales del dueño.

Hay dos etapas necesarias que se deben realizar, la selección de una macrolocalización y, dentro de ésta, la de una microlocalización definitiva.
Pero también puede darse el caso de que para una determinada empresa, la macrolocalización resulte innecesaria, dado que ya tenga preestablecida la
localidad en donde se ubicará. La selección de la macro y microlocalización está condicionada al resultado del análisis de lo que se denomina factor de localización.


\section{Factores de localización}

\img{images/12/factores.png}{Factores a considerar en un estudio de localización}

Los factores que más comúnmente influyen en la decisión de la localización de
un proyecto, se describen a continuación, aunque no se presentan en orden de
importancia y pueden existir algunos otros con base en las necesidades de cada
empresa:

\begin{enumerate}
    
\item Medios y costos de transporte. Es necesario considerar tanto el peso
como el volumen de los materiales, ya que normalmente se aplica la
tarifa que por un factor u otro resulte más alta. Además, las materias
primas, por lo general, pagan menos tarifas de transportes que los
productos terminados.
\item Disponibilidad y costo de mano de obra. Determinar cualitativa y
cuantitativamente los diversos tipos de mano de obra necesarios en la
operación de la futura planta. Además de investigar cuales son los
niveles de sueldos y salarios en las posibles localizaciones del proyecto y
su disponibilidad.
\item Localización y disponibilidad de las fuentes de abastecimiento
(materias primas). La ubicación de las materias primas resulta un factor
fundamental, ya que en ocasiones la ubicación de ciertos proyectos la
determina la fuente de materias primas. Es vital considerar los
requerimientos de insumos, condiciones de abastecimiento, costos de
transporte, existencia actual y a largo plazo, si la disponibilidad es
constante o estacional, etc.
\item Cercanía del mercado. El proyecto para la nueva ubicación debe
considerar cierta distancia entre los clientes potenciales y la localización
de la planta, ya que en muchos casos, conforme incremente la distancia,
incrementan también los costos de transportación, lo que puede
repercutir directamente en las utilidades de la empresa.
\item Factores ambientales. Este factor se debe tomar en cuenta para
aquellas zonas o lugares con climas extremos que impidan el buen
desempeño del personal o el proceso de producción de los bienes.
\item Eliminación de desechos. Para algunas plantas industriales, la
disponibilidad de medios naturales para deshacerse de ciertos desechos
resulta indispensable, por lo que su localización queda sujeta a esta
opción. En otras zonas, los reglamentos locales y gubernamentales
limitan o regulan la cantidad o la naturaleza de los desechos que pueden
arrojarse a la atmósfera o a corrientes y lechos acuosos. 

\item Costo, disponibilidad y topografía de terrenos. Es importante
considerar las necesidades actuales y las expectativas de crecimiento
futuro que pueda tener la empresa, para no tener problemas por falta de
espacio o por factores no considerados como zonas sísmicas, terrenos
extremadamente húmedos, etc. El costo del terreno puede o no ser factor
de importancia, pero dependerá del poder adquisitivo de cada empresa y
del tipo de proyecto a realizar.
\item  Estructura impositiva y legal. Algunos países o estados adoptan
políticas que promueven la instalación industrial en determinadas zonas y
ciudades creando al mismo tiempo parques industriales y ofreciendo
incentivos fiscales o de otro tipo. Conocer las leyes de cada lugar dan
un marco de restricciones y oportunidades que pueden ser de
importancia a la hora de seleccionar un lugar.
\item Disponibilidad de agua, energía eléctrica y otros suministros. La
disponibilidad tanto de agua como de energía eléctrica suelen ser un
factor determinante en la localización industrial, ya que la mayor parte de
equipos industriales modernos y sus procesos, utilizan dichos recursos y,
aunque ambos pueden ser transportables, la inversión en este tipo de
obras no se justifica para una sola empresa.
\item  Comunicaciones. Actualmente los medios de comunicación son muy
importantes, ya que estos permiten un mejor intercambio de información
entre proveedores, productores y clientes, además de disminuir costos.
\item Condiciones sociales y culturales (condiciones de vida). Se estudian
algunas variables demográficas como tamaño, distribución, edad y
cambios migratorios, además de aspectos como la actitud hacia la nueva
industria, disponibilidad, calidad y confiabilidad en los trabajadores en
potencia, tradiciones y costumbres que pueden inferir con las
modalidades conocidas de realizar negocios.
\item Disponibilidad y confiabilidad de los sistemas de apoyo. Para
algunas organizaciones, puede ser de importancia el contar con buenos
servicios públicos como facilidades habitacionales y educacionales,
caminos-vías de acceso y calles, seguridad pública, ambulancias,
servicios médicos, bomberos, etc. 


\end{enumerate}






\actividad{Cuestionario}
\begin{enumerate}
    \item ¿Qué es un estudio de localización y cuál es su objetivo principal?
    \item ¿Qué factores técnicos deben considerarse al realizar un estudio de localización?
    \item ¿Qué aspectos económicos influyen en la elección de una localización?
    \item ¿Qué rol cumplen los factores sociales y culturales en la localización de un proyecto?
    \item ¿Por qué se dice que una localización puede ser óptima en un momento y no en otro?
    \item ¿Qué importancia tienen los aspectos legales e impositivos en la localización?
    \item ¿Cómo influye la disponibilidad de infraestructura (agua, energía, transporte) en la decisión de localización?
    \item ¿Qué diferencias existen entre criterios objetivos y subjetivos al elegir una localización?
    \item ¿Qué riesgos conlleva una mala decisión de localización para una empresa?
    \item ¿Qué métodos o herramientas se utilizan para evaluar alternativas de localización?
\end{enumerate}

\actividad{Actividades para proyecto}
\begin{enumerate}
    \item Analizar un caso práctico: Seleccionar una empresa local y describir qué factores incidieron en su elección de localización.
    \item Elaborar un cuadro comparativo con al menos tres posibles localizaciones para un proyecto hipotético, evaluando sus ventajas y desventajas.
    \item Diseñar un mapa de localización (real o esquemático) que indique la ubicación más adecuada para una industria determinada.
    \item Identificar y clasificar los factores de localización en: técnicos, económicos, sociales y legales.
    \item Resolver en grupo un ejercicio de toma de decisión de localización utilizando un sistema de puntuación por criterios.
\end{enumerate}






\end{document}